\sect{11. Equipment and adventure items}

The pictures show one example from each family; variants may share an icon. The colour of a potion or ring is not enough to identify its effect.

Each item has a place of use: hand, body, inventory or dungeon mechanism. The name alone does not establish its properties. Picking up a scroll with the left mouse button shows its spell name without consuming it. Unidentified potions and wands retain generic names: use Read Magic to identify them.

\textbf{Bonuses and curses.} A positive bonus can improve attack, damage or armour. A curse can worsen equipment and prevent removal. Save before testing suspicious items. There is no camp option for removing curses. Creatures which rust or destroy equipment can cost you items: carry replacements.

\refspace
\subsection*{Weapons and ammunition}

Base damage against small/medium and large enemies, respectively, before strength, bonuses and defences. \textbf{F}: fighter, ranger and paladin; \textbf{M}: mage; \textbf{C}: cleric; \textbf{T}: thief. Multiclass combinations can broaden the options; the game checks whether the character can use the equipment.

\begin{reference}{Name}{Effect, use and limits}
\itemname{axe}{Axe} & 1d8 / 1d10. One hand; F, T. Includes Drow Cleaver.\\[4pt]
\itemname{long-sword}{Long Sword} & 1d8 / 1d12. One hand; F, T. Variants include Slasher, Night Stalker, Severious and the adamantite sword.\\[4pt]
\itemname{short-sword}{Short Sword} & 1d6 / 1d8. One hand; F, T. Slicer belongs to this family.\\[4pt]
\itemname{dagger}{Dagger} & 1d4 / 1d3. F, M, T. Thrown when used from a hand; retrieve it afterwards. Backstabber, Flicka and Guinsoo are variants.\\[4pt]
\itemname{spear}{Spear} & 1d6 / 1d8. F, T. Used as a throwing weapon from a hand. It is not replaced unless you carry another.\\[4pt]
\itemname{halberd}{Halberd} & 1d10 / 2d6. Two hands; F, T. The Chieftain Halberd is a variant of this family.\\[4pt]
\itemname{mace}{Mace / Scepter} & 1d6 + 1 / 1d6. One hand; F, C, T. The Scepter of Kingly Might functions here as a weapon of this family; do not assume a special pushing power.\\[4pt]
\itemname{flail}{Flail} & 1d6 + 1 / 2d4. One hand; F, C, T. Used to strike from the front rank.\\[4pt]
\itemname{staff}{Staff} & 1d6 / 1d6. Two hands; all classes. A mage can wield it too.\\[4pt]
\itemname{bow}{Bow} & 1d6 / 1d6. Two hands; F, T. Needs arrows in the other hand and can reload from the quiver. Includes the drow bow.\\[4pt]
\itemname{sling}{Sling} & 1d4 / 1d4. Two hands; F, C, T. Uses rocks as ammunition. Check both hands if the no-ammunition warning appears.\\[4pt]
\itemname{dart}{Dart} & 1d3 / 1d2. F, M, T. A throwing weapon; can reload from the belt. Special darts have their own modifiers.\\[4pt]
\itemname{arrow}{Arrow and Rock} & Ammunition. Arrows go in the quiver and supply the bow. Rocks supply the sling and can also be thrown. Retrieve these items when they remain on the floor.\\[4pt]
\end{reference}

\refspace
\subsection*{Armour and protection}

Lower AC is better. Base values below exclude dexterity and magical items. Two differently named items can produce the same final AC.

\begin{reference}{Name}{Effect, use and limits}
\itemname{leather}{Leather Armour} & Base AC 8; F, C, T. An early choice for a character unable to wear heavy armour.\\[4pt]
\itemname{scale}{Scale Mail} & Base AC 6; F, C. Magical bonuses and dexterity may improve the result.\\[4pt]
\itemname{chain}{Chain Mail} & Base AC 5; F, C. Equip it on the body rather than in a hand.\\[4pt]
\itemname{banded}{Banded Armour} & Base AC 4; F, C. Compare the final result on the character sheet.\\[4pt]
\itemname{plate}{Plate Mail} & Base AC 3; F, C. Includes magical or cursed variants such as Plate Mail of Great Beauty; the name does not guarantee an improvement.\\[4pt]
\itemname{robe}{Robe} & All classes. A plain robe does not itself improve AC; the Robe of Defense has its own bonus.\\[4pt]
\itemname{shield}{Shield} & Reduces AC by 1 plus any magical bonus. Must occupy a usable hand; it does not combine with using a two-handed weapon.\\[4pt]
\itemname{bracers}{Bracers} & Their bonus may improve AC when worn on the wrists. Compare the final value after changing equipment.\\[4pt]
\itemname{ring-protection}{Ring of Protection} & Wear it on a finger. The better of two rings is used, and it does not add to magically enhanced armour.\\[4pt]
\itemname{helmet}{Helmet and Boots} & Equip them in the head and feet slots. Ordinary helmets and boots do not improve AC.\\[4pt]
\end{reference}

\refspace
\subsection*{Tools, keys and special items}

\begin{reference}{Name}{Effect, use and limits}
\itemname{spellbook}{Spell Book} & A mage needs it to prepare magic. Put it in a hand and right-click to open prepared spells. Casting does not consume the book.\\[4pt]
\itemname{holy-symbol}{Holy Symbol} & A cleric or paladin's tool. In a hand it opens prayers. Carrying it does not replace selecting and recovering spells through rest.\\[4pt]
\itemname{lockpicks}{Lockpicks} & Pick them up with the pointer and use them on a suitable lock. A thief can attempt to open it; failure does not necessarily mean the object is not a lock.\\[4pt]
\itemname{key}{Keys} & There are gold, silver, ruby, jeweled, skull, dwarven and drow keys. Bring the key to the matching slot or lock. A similar colour or shape does not guarantee a match.\\[4pt]
\itemname{gem}{Gems} & Keep gems: they may serve as offerings or mechanism pieces. Their use depends on the location; they are not currency to spend in a shop.\\[4pt]
\itemname{stone-object}{Stone Objects} & Stone dagger, holy symbol, medallion, necklace, orb, ring and scepter. These are travel-mechanism pieces. Their shape matters; keep each one and note where it was found.\\[4pt]
\itemname{bones}{Bones} & Human, dwarf and halfling remains can be connected to characters in the adventure. Keep them and pay attention to dialogue; they are not food.\\[4pt]
\itemname{letter}{Commission Letter and Notes} & Put a document in a hand and use it to read it. Notes provide context or clues. They are documents, not an automap that updates as you walk.\\[4pt]
\itemname{orb}{Orb of Power} & A special adventure item. Do not treat it as a rechargeable wand or ordinary weapon; keep unique objects until their purpose becomes clear.\\[4pt]
\itemname{necklace}{Necklaces and Medallions} & Worn on the neck. Some are ornaments, including the Luck Stone Medallion, which gives no combat bonus.\\[4pt]
\end{reference}
